\documentclass[10pt]{article}
\usepackage[margin=1in]{geometry}
\usepackage[T1]{fontenc}
\usepackage{lmodern,microtype,amsmath,amssymb,amsthm}
\theoremstyle{definition}
\newtheorem{assumption}{Assumption}
\begin{document}
\section*{Stochastic Generalized Lipschitz Convex Minimization}
We consider minimizing a continuous convex function $f:E\to\mathbb{R}$ on a separable real Hilbert space $E$. A given point $x^0\in E$ satisfies $\|x^0-x^\star\|\leq R$, where $R>0$ and $x^\star$ is a global minimizer. At $x\in E$, the oracle returns a random vector whose mean is a subgradient of $f$. Each reply decomposes into a component with second moment at most $G_0^2$ and a component bounded almost surely in norm by $G_1(f(x)-f(x^\star))$, where $G_0>0$ and $G_1\geq0$. Samples are independent and identically distributed. For $\varepsilon>0$, the goal is a measurable output $\bar{x}\in E$ satisfying $f(\bar{x})-f(x^\star)\leq\varepsilon$ with probability at least $9/10$.

The optimization problem is
\[
\min_{x\in E} f(x).
\]
Write $\langle\cdot,\cdot\rangle$ for the inner product and $\|\cdot\|$ for its induced norm. The algorithm receives $E$, $x^0$, $R$, $G_0$, $G_1$, and $\varepsilon$; the objective and sampling distribution are accessed through the oracle.

Our assumptions:
\begin{assumption}[Convexity and initial distance]
The function $f$ is continuous, and for every $x,y\in E$ and $t\in[0,1]$,
\[
f((1-t)x+ty)\leq (1-t)f(x)+tf(y).
\]
There exists a fixed $x^\star\in E$ such that
\[
\begin{aligned}
f(x^\star)&\leq f(x) &&\text{for every }x\in E,\\
\|x^0-x^\star\|&\leq R.
\end{aligned}
\]
\end{assumption}

\begin{assumption}[Stochastic subgradient oracle]
Let $\Omega$ be a measurable space, let $\nu$ be a probability measure on $\Omega$, and let $g:E\times\Omega\to E$ be jointly measurable. For every $x\in E$, the map $g(x,\cdot)$ is Bochner integrable under $\nu$, and
\[
f(x)-f(y)\leq
\left\langle\mathbb{E}_{\xi\sim\nu}[g(x,\xi)],x-y\right\rangle,
\qquad x,y\in E.
\]
At call $k=0,1,\ldots$, a query $x\in E$ receives $g(x,\xi_k)$. The random variables $(\xi_k)_{k\geq0}$ are measurable, mutually independent, and each has law $\nu$. Algorithm randomness is independent of this sample sequence.
\end{assumption}

\begin{assumption}[Generalized moment bound]
For every $x\in E$, there exist measurable maps $u_x,v_x:\Omega\to E$ such that $\|u_x(\cdot)\|^2$ is integrable under $\nu$ and
\[
\begin{aligned}
g(x,\xi)&=u_x(\xi)+v_x(\xi)
&&\nu\text{-a.s.},\\
\mathbb{E}_{\xi\sim\nu}\bigl[\|u_x(\xi)\|^2\bigr]&\leq G_0^2,\\
\|v_x(\xi)\|&\leq G_1\bigl(f(x)-f(x^\star)\bigr)
&&\nu\text{-a.s.}
\end{aligned}
\]
\end{assumption}

After the allotted oracle calls, the algorithm must return a measurable output $\bar{x}\in E$ satisfying
\[
P\bigl(f(\bar{x})-f(x^\star)>\varepsilon\bigr)\leq\frac{1}{10},
\]
where $P$ is the joint law of the algorithm's randomness and the oracle sample sequence.
\end{document}
